\section{The ecosystems did not run the same experiment}
\label{sec:implementation}

The manuscript states a common protocol: the same 30 epochs, learning rate,
optimiser and batch size. Reading the source of all eight implementations shows
otherwise.

\begin{table}[t]
\centering
\caption{Training protocol as implemented, read from source. The manuscript
states a common learning rate; two values are in use.}
\label{tab:protocol}
\small
\begin{tabular}{llrlrl}
\toprule
Ecosystem & Optimiser & LR & Loader mechanism & Threads & Input scaling \\
\midrule
Python/PyTorch    & Adam & 1e-4 & DataLoader workers & 2 & $[0,1]$ \\
Python/TensorFlow & Adam & \textbf{1e-3} & Keras generator & 1 & $[0,1]$ \\
Python/JAX        & Adam & 1e-4 & Keras generator & 1 & $[0,1]$ \\
C++/LibTorch      & Adam & 1e-4 & DataLoader workers & 2 & $[0,1]$ \\
Java/DL4J         & Adam & 1e-4 & AsyncDataSetIterator & 2 & $[0,1]$ \\
R/torch           & Adam & 1e-4 & dataloader workers & \textbf{0} & $[0,1]$ \\
Rust/tch          & Adam & \textbf{1e-3} & rayon, all cores & \textbf{96} & \textbf{mean/std} \\
MATLAB/DLT        & adam & 1e-4 & augmentedImageDatastore & 1 & $[0,1]$ \\
\bottomrule
\end{tabular}
\end{table}

\subsection{Data-loader parallelism, not language, tracks the ranking}

The largest single difference is how many CPU threads decode and feed images: from
zero (R decodes on the main thread) to all 96 cores (Rust, through \code{rayon}).
Spearman correlation between loader threads and mean epoch duration is
\(\rho=-0.73\) (\(p=0.04\)) across the eight ecosystems. The fastest stack
decodes on 96 cores; the slowest uses no loader workers; they differ by
\(11.6\times\) in duration.

Combined with the GPU-load result of Section~\ref{sec:instrument} --- the
accelerator at a fraction of its limit, duration accounting for essentially the
whole log energy spread --- the most parsimonious reading of the headline result
is that it ranks data-pipeline configurations.

\subsection{One ecosystem was not solving the same task}

Three defects in the Rust implementation all push in the direction of the
reported result:

\begin{itemize}\itemsep2pt
  \item \textbf{Fashion-MNIST was trained at \(28\times28\times1\).} The loader
        applied no resize and no channel replication, so training ran on the
        native grayscale image while \emph{its own evaluation loop} --- and every
        other ecosystem, in both phases --- used \(32\times32\times3\). The
        \code{conv1} input volume during training was \(0.26\times\) everyone
        else's.
  \item \textbf{Tiny ImageNet was trained at its native \(64\times64\)} in the
        ResNet binary, while the VGG binary passed \code{Some(32)} and every
        other ecosystem used \(32\times32\): four times the spatial work for the
        same nominal configuration, inside one ecosystem.
  \item \textbf{Evaluation ran one image at a time} (\code{iter\_batches(1)}) in
        all six binaries, against batch 128 elsewhere. Batch-1 GPU inference is
        launch-overhead bound.
\end{itemize}

The published numbers line up with the defects. Against the median of the other
stacks, this ecosystem's training energy is \(0.07\times\) for
ResNet-18/Fashion-MNIST --- the cheapest cell in the entire campaign, and
precisely the cell with the smallest input --- rising to \(0.61\times\); its
inference energy is \(1.51\)--\(1.95\times\) on CIFAR-100 and Fashion-MNIST.

The manuscript's \emph{training/inference ranking reversal}, reported as a
finding, is therefore reproducible for this ecosystem as a pure artefact of input
shape and evaluation batching.

\subsection{One ecosystem evaluated in training mode}

The TensorFlow stack built its functional graph with
\code{backbone(inputs, training=True)}, pinning training-mode behaviour into the
graph. Batch normalisation then used the statistics of the current batch during
evaluation as well. Because the test split is served in class order, each
evaluation batch held a single class and normalisation ran against degenerate
per-class statistics: measured test accuracy was \SI{21}{\percent} against
\SI{83}{\percent} on train, and \SI{86}{\percent} once the flag was removed.

Two consequences. The accuracy this stack would have reported was meaningless ---
but accuracy was never written to disk, so nothing surfaced it. And its inference
phase was measuring a different computation from every other ecosystem, all of
which switch to evaluation mode.

\subsection{The dataset itself was ambiguous}

Fashion-MNIST class~0 is \code{T-shirt/top}. Used verbatim as a directory name it
creates a \emph{nested} directory, so \code{train/T-shirt/} holds a \code{top/}
subdirectory rather than images. Loaders then disagree about what the dataset
contains: \code{ImageFolder} and \code{flow\_from\_directory} walk recursively and
recover the \num{6000} images, a loader listing only direct children sees an empty
class, and a loader using \code{read\_dir} yields the subdirectory as an image
path. The same dataset directory is three different datasets depending on who
reads it.
